\section*{0.1 The Goal: Quantifying Topological Knowledge Drift}
\addcontentsline{toc}{section}{0.1 The Goal: Quantifying Topological Knowledge Drift}

In classical statistical mechanics, a system's macroscopic state is derived by integrating over the phase space of its microscopic constituents. We apply an isomorphic principle to academic literature. Rather than tracking particles in $\mathbb{R}^3$, we map natural language into a high-dimensional vector space $\mathbb{R}^d$ ($d=768$ for standard Transformer architectures).

Let this embedding space represent a non-linear Riemannian manifold $\mathcal{M}$, where the distance metric between two vectors is not the standard Euclidean $L_2$ norm, but a cosine distance derived from the attention mechanism's learned probability distributions.

A single academic concept moving through this space over time $t$ can be modeled as a stochastic trajectory—analogous to a Wiener process. We define \textbf{Topological Knowledge Drift} as the measurable derivative of lexical and semantic clustering across temporal epochs. By applying band-pass signal filtering to the frequency derivatives $\frac{\partial f}{\partial t}$ of localized jargon, we can mathematically isolate the genesis of new scientific sub-disciplines from the ambient noise of general academic boilerplate.
